% Fragment showing expected style for a step-by-step derivation.
% Adapt content to the blueprint; do not copy unless the topic matches.

\subsection{Deriving the computational formula for variance}

Let \(X\) be a random variable with finite second moment, and let
\(\mu = \mathbb{E}[X]\). The goal is to rewrite the definition of variance
in a form that is often easier to compute.

\begin{align}
  \operatorname{Var}(X)
    &= \mathbb{E}\left[(X-\mu)^2\right] \\
    &= \mathbb{E}\left[X^2 - 2\mu X + \mu^2\right] \\
    &= \mathbb{E}[X^2] - 2\mu \mathbb{E}[X] + \mu^2 \\
    &= \mathbb{E}[X^2] - 2\mu^2 + \mu^2 \\
    &= \mathbb{E}[X^2] - \mu^2.
\end{align}

The third line uses linearity of expectation and the fact that \(\mu\) is a
constant. Since \(\mu=\mathbb{E}[X]\), the final form is
\[
  \boxed{\operatorname{Var}(X)=\mathbb{E}[X^2]-(\mathbb{E}[X])^2.}
\]

As a quick check, if \(X=c\) is constant, then both sides give \(0\), as they
should: a constant random variable has no spread.
